\section{Stability}\label{sec:stability}

This section proves that a moving sofa whose area is close to $\abs\Gs$ is,
after a translation, close to $\Gs$, at the rate of the square root of the
missing area, and that this rate is optimal. As in \cref{sec:equality},
$\varphi$ is Gerver's angle, $\KG$ is Gerver's cap and $\cA=\cA_{\pi/2}$. The
\emph{deficit} of a set $S$ of finite area is $\eps(S)=\abs\Gs-\abs S$, which
is nonnegative for moving sofas (\cref{fact:optimal}). Two sets
$X,Y\subset\R^2$ are \emph{$r$-close} if every point of each lies within
Euclidean distance $r$ of a point of the other; for nonempty compact sets, this
means that their Hausdorff distance is at most $r$. For a nonempty compact set
$S$, let $h_S(t)=\max_{p\in S}\ip p{u_t}$. The \emph{normalized} translate
$S^\natural=S+(h_S(\pi)-h_\Gs(\pi),1-h_S(\pi/2))$ has
$h_{S^\natural}(\pi/2)=1$ and $h_{S^\natural}(\pi)=h_\Gs(\pi)$, and $S$ is
\emph{normalized} if $S^\natural=S$. Let $X\triangle Y=(X\setminus
Y)\cup(Y\setminus X)$ be the symmetric difference.

\begin{theorem}[stability of Gerver's sofa]\label{thm:stability}
There are $C,C',\eps_0>0$ such that for every moving sofa $S$ with
$\eps=\eps(S)<\eps_0$, the sets $S^\natural$ and $\Gs$ are
$C\sqrt\eps$-close, and $\abs{S^\natural\triangle\Gs}\le C'\sqrt\eps$.
\end{theorem}

\begin{theorem}[the rotation angle]\label{thm:stability-angle}
There are $C'',\eps_0>0$ such that every moving sofa $S$ with rotation angle
$\omega\in[\omega_0,\pi/2]$ and $\eps(S)<\eps_0$ satisfies
$0\le\pi/2-\omega\le C''\eps(S)$.
\end{theorem}

A \emph{rigid motion} is a map $g(p)=R_\vartheta p+v$. For nonempty compact
sets $X,Y$, $d_{\mathrm{rig}}(X,Y)$ is the infimum of the $d\ge0$ for which
$X$ and $g(Y)$ are $d$-close for some rigid motion $g$. Let $\mathbb B(p,r)$
and $\mathbb B^\circ(p,r)$ be the closed and the open disk of radius $r$ about
$p$.

\begin{theorem}[the exponent $1/2$ is optimal]\label{thm:stability-exponent}
\begin{enumerate}[label=(\alph*)]
\item\label{thm:stability-exponent:family} There are a point $p$ and $r_0>0$
such that for every $r\in(0,r_0)$ the set $\Gs_r=\Gs\setminus\mathbb
B^\circ(p,r)$ is a moving sofa with $\eps(\Gs_r)=\pi r^2$ and
$d_{\mathrm{rig}}(\Gs_r,\Gs)=r$. Moreover, $\Gs_r$ and $\Gs$ are $r$-close,
and $r\le d$ whenever $\Gs_r$ and $g(\Gs)$ are $d$-close for a rigid motion
$g$.
\item\label{thm:stability-exponent:no} Let $a>1/2$, $C\in\R$ and $\eps_0>0$.
There is a moving sofa $S$ with $0<\eps(S)<\eps_0$ such that $S$ and $g(\Gs)$
are $C\eps(S)^a$-close for no rigid motion $g$.
\item\label{thm:stability-exponent:rig} For the same $a$, $C$ and $\eps_0$,
there is a moving sofa $S'$ with $0<\eps(S')<\eps_0$ and
$C\eps(S')^a<d_{\mathrm{rig}}(S',\Gs)$.
\end{enumerate}
\end{theorem}

Baek's functional $\cQ$ is given by a formula \baek{Def.~8.2.2}: for convex
bodies $K$, $B$ and $D$,
\begin{multline}\label{eq:stab-Q}
  \cQ(K,B,D)=\abs K
  +\tfrac12\int_{(3\pi/2,\,3\pi/2+\varphi^{\mathrm L})}h_D\dd\sigma_D
  +\mathcal J\bigl([Y_D,\xx_K(\varphi^{\mathrm L})]\bigr)\\
  -\mathcal J\bigl(\xx_K|_{[\varphi^{\mathrm R},\varphi^{\mathrm L}]}\bigr)
  +\mathcal J\bigl([\xx_K(\varphi^{\mathrm R}),X_B]\bigr)
  +\tfrac12\int_{(\pi+\varphi^{\mathrm R},\,3\pi/2)}h_B\dd\sigma_B ,
\end{multline}
where $X_B=v_B^+(\pi+\varphi^{\mathrm R})$,
$Y_D=v_D^-(3\pi/2+\varphi^{\mathrm L})$ and $\mathcal J([p,q])=\frac12\,p\times
q$; the integrals are the curve areas of two arcs of the boundaries of $D$ and
$B$ \baek{Thm.~7.3.2}. Let $\bar\cT$ be defined as $\cT$ (\cref{sec:Q}), with
$K\in\Kc_{\pi/2}$ in place of $K\in\Ki$; by \cref{lem:stab-variation} below,
$\cQ\le\abs\Gs$ on $\bar\cT$. For $K\in\Kc_{\pi/2}$ let
$s_K=h_{\KG}(\pi)-h_K(\pi)$, so that $K$ and $\KG+(s_K,0)$ have the same
leftmost abscissa.

\begin{theorem}[the cap estimate]\label{thm:stability-cap}
\begin{enumerate}[label=(\alph*)]
\item\label{thm:stability-cap:wide} For every $\xi=(K,B,D)\in\bar\cT$, the
caps $K$ and $\KG+(s_K,0)$ are $2\sec\varphi\sqrt{\abs\Gs-\cQ(\xi)}$-close.
\item\label{thm:stability-cap:ki} For every $K\in\Ki$, the caps $K$ and
$\KG+(s_K,0)$ are $2\sec\varphi\sqrt{\abs\Gs-\cA(K)}$-close.
\item\label{thm:stability-cap:num} $2\sec\varphi\le2500/1249<1001/500$, so
\ref{thm:stability-cap:wide} and \ref{thm:stability-cap:ki} hold with
$1001/500$ in place of $2\sec\varphi$.
\end{enumerate}
\end{theorem}

\Cref{sec:stab-certificate,sec:stab-cap} bound, for every
$\xi=(K,B,D)\in\bar\cT$, the distance from $K$ to $\KG+(s_K,0)$ by a multiple
of $\sqrt{\abs\Gs-\cQ(\xi)}$: with the constant $80$ in
\cref{prop:stab-coercive}, which the proofs of
\cref{thm:stability,thm:stability-angle} use, and with $2\sec\varphi$ in
\cref{thm:stability-cap}. Near $\KG$, Baek's bound $\cA(K)\le\cQ(\xi_K)$ holds
without the injectivity condition (\cref{sec:stab-local}), and a rotation
angle $\omega<\pi/2$ costs area proportional to $\pi/2-\omega$
(\cref{sec:stab-angle}). \Cref{sec:stab-recovery} passes from the cap back to
the sofa, and a compactness argument brings the normalized translate of every
moving sofa of small deficit near $\Gs$ (\cref{sec:stab-entry}); this proves
\cref{thm:stability,thm:stability-angle}. \Cref{sec:stab-puncture} removes
small disks from $\Gs$ to prove \cref{thm:stability-exponent}.

\subsection{The deficit of \texorpdfstring{$\cQ$}{Q} on
\texorpdfstring{$\bar\cT$}{the enlarged domain}}\label{sec:stab-certificate}

Let $\mathcal D$ be a set closed under the Minkowski combinations $c_\lambda$
(\cref{sec:Q}). A functional $\Psi$ on $\mathcal D$ is \emph{quadratic} if
$\Psi(x)=\beta(x,x)$ with $\beta$ convex-linear in each argument; its
\emph{directional derivative} $D\Psi(x;y)$ is the right derivative of
$\lambda\mapsto\Psi(c_\lambda(x,y))$ at $0$ \baek{Def.~7.1.4, Def.~7.1.5}.
Baek shows that $D\Psi(x;y)=\beta(x,y)+\beta(y,x)-2\beta(x,x)$, and that a
quadratic concave $\Psi$ is largest at $x$ if and only if $D\Psi(x;y)\le0$ for
all $y$ \baek{Lemma~7.1.4, Thm.~7.1.5}. Let
$E_\Psi(x,y)=4\Psi(c_{1/2}(x,y))-2\Psi(x)-2\Psi(y)$. As $\beta$ is
convex-linear in each argument,
$E_\Psi(x,y)=\beta(x,y)+\beta(y,x)-\beta(x,x)-\beta(y,y)$, and Baek's formula
for $D\Psi$ gives
\begin{equation}\label{eq:stab-segment}
  \Psi(x)-\Psi(y)=-D\Psi(x;y)+E_\Psi(x,y).
\end{equation}
For a convex body $K$ and a curve $\mathbf z$ as in \cref{sec:Q}, we call
$\varrho_{K,\mathbf z}$ the \emph{displacement} of $\mathbf z$. For convex
bodies $K_0,K_1$, the \emph{energy} $E_\cS(K_0,K_1)$ is the sum, over the four
terms $\cM_K(a,b;\mathbf z_K)$ of $\cS_K$, of
$\frac12\int_a^b(\varrho_{K_0,\mathbf z_{K_0}}-\varrho_{K_1,\mathbf
z_{K_1}})^2\dd t$, and the energies $E_\cR(B_0,B_1)$ and $E_\cL(D_0,D_1)$ are
defined likewise. As the curves of these terms are convex-linear in the body
(\cref{fact:Q}\ref{fact:Q:decomp}), \eqref{eq:gap} gives, for all convex bodies
and $\lambda\in[0,1]$,
\begin{equation}\label{eq:stab-gap}
  (1-\lambda)\cS_{K_0}+\lambda\cS_{K_1}-\cS_{(1-\lambda)K_0+\lambda K_1}
    =\lambda(1-\lambda)E_\cS(K_0,K_1),
\end{equation}
and likewise for $\cR$ and $\cL$, which are therefore convex. Let
$I=[\varphi,\pi/2-\varphi]$.

\begin{lemma}[$\cQ$ on $\bar\cT$]\label{lem:stab-wide}
$\bar\cT$ is closed under Minkowski combinations. On $\bar\cT$,
$\cQ(K,B,D)=\Lambda(K)-\cS_K-\cR_B-\cL_D$, where $\Lambda$, Baek's
$\cP_K+\cS_K$ (\cref{fact:Q}\ref{fact:Q:decomp}), is convex-linear on
$\Kc_{\pi/2}$, and $\cQ$ is quadratic and concave.
\end{lemma}

\begin{proof}
Support functions of Minkowski combinations are the same combinations of the
support functions, so $\bar\cT$ is closed under them. Baek's proof of the
decomposition in \cref{fact:Q}\ref{fact:Q:decomp} \baek{Lemma~8.3.4} uses only
the equalities at the ends of the intervals, so it holds on $\bar\cT$. On
$\Kc_{\pi/2}$, $\abs K$ is the sum of the curve areas of the four arcs of
$\cS_K$ and of the areas $\frac12\sigma_K(t)h_K(t)$ of the triangles spanned by
$O$ and the edges $e_K(t)$, $t=0,\varphi,\pi/2-\varphi,\pi/2,\pi$. Each edge
lies on the supporting line that carries the adjacent end segments
$[v_K^+(a),\mathbf z(a)]$ and $[\mathbf z(b),v_K^-(b)]$ of the Mamikon terms.
In $\cS_K+\cP_K$ \baek{Def.~8.3.4}, the curve areas of the arcs cancel against
$\abs K$, and what remains are $\mathcal J(\yy_K|_I)$, $-\mathcal J(\xx_K|_I)$,
the triangles of the edges and curve areas of segments on the five lines
$l_K(t)$, $t=0,\varphi,\pi/2-\varphi,\pi/2,\pi$. Joining the collinear pieces
on each of these lines gives, with $W_K$ and $Z_K$ as in
\cref{sec:pinned-bounds},
\begin{multline*}
  \Lambda(K)=h_K(0)+h_K(\pi)
  +\frac{2\sin\varphi-h_K(\varphi)-h_K(\pi-\varphi)}{2\cos\varphi}\\
  +\mathcal J(\yy_K|_I)-\mathcal J(\xx_K|_I)+\varsigma_K ,
\end{multline*}
where
\begin{multline*}
  \varsigma_K=\mathcal J([\mathbf l^{\pi/2}_K(\varphi),\yy_K(\varphi)])
    -\mathcal J([W_K(\varphi),\xx_K(\varphi)])\\
  -\mathcal J([\mathbf l^{\pi-\varphi}_K(\pi/2),\yy_K(\pi/2-\varphi)])
    +\mathcal J([Z_K(\pi/2-\varphi),\xx_K(\pi/2-\varphi)]).
\end{multline*}
The differences $\yy_K-\xx_K=u_t+v_t$,
$\mathbf l^{\pi/2}_K(\varphi)-W_K(\varphi)=((1-\sin\varphi)/\cos\varphi,1)$
and
$\mathbf l^{\pi-\varphi}_K(\pi/2)-Z_K(\pi/2-\varphi)
=(-(1-\sin\varphi)/\cos\varphi,1)$ do not depend on $K$, so $\Lambda$ is
convex-linear \baek{Lemma~7.1.6}. The terms of \eqref{eq:stab-Q} are quadratic
\baek{Thm.~7.1.3, Thm.~7.3.2, Thm.~7.1.2, Prop.~7.2.2}, so $\cQ$ is quadratic;
as $\Lambda$ is convex-linear and $\cS,\cR,\cL$ are convex by
\eqref{eq:stab-gap}, $\cQ=\Lambda-\cS-\cR-\cL$ is concave.
\end{proof}

\begin{lemma}[the directional derivative at Gerver's
triple]\label{lem:stab-variation}
For every $\xi\in\bar\cT$, $D\cQ(\xi_G;\xi)\le0$. Hence
$\cQ\le\cQ(\xi_G)=\abs\Gs$ on $\bar\cT$.
\end{lemma}

\begin{proof}
Let $\xi_G=(K_0,B_0,D_0)$ and $\xi=(K,B,D)$, let $\breve\sigma_X$ be the image
of $\sigma_X$ under $t\mapsto t-\pi$, and let $\Delta=h_K-h_{K_0}$,
$\Delta_B(t)=h_B(t+\pi)-h_{B_0}(t+\pi)$ and
$\Delta_D(t)=h_D(t+\pi)-h_{D_0}(t+\pi)$. Baek computes $D\cQ(\xi_G;\xi)$ for
$\xi\in\cT$ \baek{Thm.~8.5.6, Thm.~8.5.7}, using $X_{B_0}=\xx_{K_0}(\varphi)$
and $Y_{D_0}=\xx_{K_0}(\varphi^{\mathrm L})$ \baek{Thm.~8.4.3}. The computation
differentiates only the curves of $\xi_G$, so it holds on $\bar\cT$, and it
gives
\begin{multline*}
  D\cQ(\xi_G;\xi)=\int_{[0,\pi]}\Delta\dd\sigma_{K_0}
  -\int_{I\cup(I+\pi/2)}\Delta\,i_{K_0}\dd t\\
  +\int_{(\varphi,\pi/2)}\Delta_B\dd\breve\sigma_{B_0}
  +\int_{(\pi/2,\pi-\varphi)}\Delta_D\dd\breve\sigma_{D_0}.
\end{multline*}
The four terms are the directional derivatives of $\abs K$, of $-\mathcal
J(\xx_K|_I)$, which is an integral of $\Delta$ against a density $i_{K_0}$, of
the terms of \eqref{eq:stab-Q} that involve $B$ and of those that involve $D$
\baek{Thm.~8.5.2, Thm.~8.5.3}. The first holds for all caps, as $\Delta$
vanishes at $3\pi/2$, the only angle in $(\pi,2\pi)$ that carries
$\sigma_{K_0}$.

Baek shows that, on $[0,\pi]$, $\sigma_{K_0}$ is the sum of the measure with
density $i_{K_0}$ on $I\cup(I+\pi/2)$, of $\breve\sigma_{B_0}$ on
$[\pi/2-\theta,\pi/2)$, of $\breve\sigma_{D_0}$ on $(\pi/2,\pi/2+\theta]$ and of
the atom at $\pi/2$ \baek{Thm.~8.4.5}, where $\theta$ is Gerver's second angle
(\cref{sec:gerver}). Moreover, $\breve\sigma_{B_0}$ and $\breve\sigma_{D_0}$
vanish on $(\varphi,\pi/2-\theta)$ and $(\pi/2+\theta,\pi-\varphi)$, and
$h_{K_0}(t)+h_{B_0}(t+\pi)=1$ on $[\pi/2-\theta,\pi/2]$ and
$h_{K_0}(t)+h_{D_0}(t+\pi)=1$ on $[\pi/2,\pi/2+\theta]$ \baek{Thm.~8.4.3}. So
the part of the first term against the density $i_{K_0}$ cancels the second
term, the atom contributes $\Delta(\pi/2)\sigma_{K_0}(\{\pi/2\})=0$ as
$h_K(\pi/2)=h_{K_0}(\pi/2)$, and the rest of the first term combines with the
last two terms into
\begin{multline*}
  D\cQ(\xi_G;\xi)=\int_{[\frac\pi2-\theta,\frac\pi2)}
    \bigl(h_K(t)+h_B(t+\pi)-1\bigr)\dd\breve\sigma_{B_0}(t)\\
  +\int_{(\frac\pi2,\frac\pi2+\theta]}
    \bigl(h_K(t)+h_D(t+\pi)-1\bigr)\dd\breve\sigma_{D_0}(t),
\end{multline*}
which is nonpositive by the inequalities defining $\bar\cT$. By
\cref{lem:stab-wide} and \baek{Thm.~7.1.5}, $\cQ\le\cQ(\xi_G)$ on $\bar\cT$,
and $\cQ(\xi_G)=\abs\Gs$ (\cref{fact:Q}\ref{fact:Q:max}).
\end{proof}

\begin{proposition}[the deficit certificate]\label{prop:stab-certificate}
For every $\xi=(K,B,D)\in\bar\cT$,
\[
  \abs\Gs-\cQ(\xi)=-D\cQ(\xi_G;\xi)+E_\cS(\KG,K)+E_\cR(B_{\KG},B)
    +E_\cL(D_{\KG},D),
\]
and the four terms on the right are nonnegative; so
$E_\cS(\KG,K)\le\abs\Gs-\cQ(\xi)$.
\end{proposition}

\begin{proof}
Apply \eqref{eq:stab-segment} to $\cQ$, $x=\xi_G$ and $y=\xi$. By
\cref{lem:stab-wide}, the convex-linear $\Lambda$ drops out of
$E_\cQ(\xi_G,\xi)$, and \eqref{eq:stab-gap} at $\lambda=1/2$ gives
$E_\cQ(\xi_G,\xi)=E_\cS(\KG,K)+E_\cR(B_{\KG},B)+E_\cL(D_{\KG},D)$. The first
term is nonnegative by \cref{lem:stab-variation}, and the energies are
integrals of squares.
\end{proof}

\subsection{From the energies to the cap; proof of
\texorpdfstring{\cref{thm:stability-cap}}{the cap estimate}}\label{sec:stab-cap}

For any function $f$ with right derivatives and any $T\in\R$, the
\emph{tangent residual} $r_T$ and the \emph{corner residual} $r_{\mathrm c}$
of $f$ are
\[
  r_T(t)=\frac{f(T)-f(t)\cos(T-t)}{\sin(T-t)}-f'(t^+),\qquad
  r_{\mathrm c}(t)=f(t+\pi/2)-f'(t^+).
\]
The \emph{four residuals} of $f$ are $r_1=r_{\pi/2}$ on $(0,\varphi)$,
$r_2=r_{\mathrm c}$ on $(\varphi,\pi/2-\varphi)$, $r_3=r_{\pi-\varphi}$ on
$(\pi/2-\varphi,\pi/2)$ and $r_4=r_\pi$ on $(\pi/2,\pi)$, and
$E(f)=\frac12\sum_j\int r_j^2$, each integral over the interval of $r_j$. By
the quotient rule, $q_T(t)=(f(t)-f(T)\cos(T-t))/\sin(T-t)$ has the right
derivative $-r_T(t)/\sin(T-t)$. So, if $f$ is continuous with integrable right
derivatives on $[s_1,s_2]$, the fundamental theorem of calculus for right
derivatives gives
\begin{gather}
  q_T(s_1)-q_T(s_2)=\int_{s_1}^{s_2}\frac{r_T(u)}{\sin(T-u)}\dd u\quad
  \text{if $\sin(T-u)\ne0$ on $[s_1,s_2]$},\label{eq:stab-ftc}\\
  f(s_1)=f(s_2)-\int_{s_1}^{s_2}f(u+\pi/2)\dd u
    +\int_{s_1}^{s_2}r_{\mathrm c}(u)\dd u.\notag
\end{gather}
For $K_0,K_1\in\Kc_{\pi/2}$ and $\Delta=h_{K_1}-h_{K_0}$, the function
\emph{formed with} $K_0$ and $K_1$ is $f(t)=\Delta(t)+\Delta(\pi)\cos t$. Then
$f=h_{K_1}-h_{K_0+(s,0)}$ with $s=h_{K_0}(\pi)-h_{K_1}(\pi)$,
$f(\pi/2)=f(\pi)=0$, and $f$ has at every $t$ the right derivative
$f'(t^+)=\ip{v^+_{K_1}(t)-v^+_{K_0}(t)}{v_t}-\Delta(\pi)\sin t$
(\cref{fact:convex}\ref{fact:convex:vertices}).

\begin{lemma}[residuals and displacements]\label{lem:stab-residuals}
Let $K_0,K_1\in\Kc_{\pi/2}$, and let $f$ be formed with them. On the open
interval of each term $\cM_K(a,b;\mathbf z_K)$ of $\cS_K$, the corresponding
residual of $f$ is $\varrho_{K_1,\mathbf z_{K_1}}-\varrho_{K_0,\mathbf
z_{K_0}}$. The four residuals and their squares are integrable, and
$E(f)=E_\cS(K_0,K_1)$.
\end{lemma}

\begin{proof}
As in the proof of \cref{lem:tangent},
$\varrho_{K,\mathbf l_K^T}(t)=(h_K(T)-h_K(t)\cos(T-t))/\sin(T-t)
-\ip{v^+_K(t)}{v_t}$ for $T-\pi<t<T$, and
$\varrho_{K,\yy_K}(t)=h_K(t+\pi/2)-\ip{v^+_K(t)}{v_t}$. So
$\varrho_{K,\mathbf l^T_K}$ and $\varrho_{K,\yy_K}$ are the residuals $r_T$ and
$r_{\mathrm c}$ of $h_K$, whose right derivative is $\ip{v^+_K(t)}{v_t}$
(\cref{fact:convex}\ref{fact:convex:vertices}). Residuals are linear in the
function, and multiples of $\cos t$, whose right derivatives are the same
multiples of $-\sin t$, have zero residuals; as
$f=h_{K_1}-h_{K_0}+\Delta(\pi)\cos t$, the residuals of $f$ are the
differences of the displacements. By Mamikon's theorem \baek{Thm.~7.4.1}, the
displacements are measurable and bounded, so the residuals and their squares
are integrable.
\end{proof}

\begin{proposition}[the support bound]\label{prop:stab-coercive}
Let $0<\varphi<\pi/4$, not necessarily Gerver's angle, and form the four
residuals with this $\varphi$. Let $f$ be continuous, with right derivatives on
$(0,\pi)$, $f(\pi/2)=f(\pi)=0$, and four residuals that are integrable and
square integrable. For $t\in[0,\pi]$:
\begin{enumerate}[label=(\alph*)]
\item\label{prop:stab-coercive:80}
$\abs{f(t)}\le20\sum_{j=1}^4\int\abs{r_j}\le80\sqrt{E(f)}$;
\item\label{prop:stab-coercive:sharp} $f(t)^2\le\mathcal G_\varphi(t)\sum_j\int
r_j^2\le2\sec^2\varphi\sum_j\int r_j^2$, so
$\abs{f(t)}\le2\sec\varphi\sqrt{E(f)}$, where $\mathcal G_\varphi(t)$ is
$\cos^2t\,(2\sec^2\varphi-\tan t)$ for $t\le\varphi$,
$\cos t\,(2\sec\varphi-\sin t)$ for $\varphi<t\le\pi/2-\varphi$,
$\sin t\cos t+2\tan\varphi\cos^2t$ for $\pi/2-\varphi<t\le\pi/2$, and
$-\sin t\cos t$ for $t>\pi/2$.
\end{enumerate}
For Gerver's angle $\varphi$, if $\xi=(K,B,D)\in\bar\cT$ and $f$ is formed with
$\KG$ and $K$, then $\abs f\le80\sqrt{\abs\Gs-\cQ(\xi)}$ on $[0,\pi]$.
\end{proposition}

\begin{proof}
Let $b'=\pi/2-\varphi$ and $T=\pi-\varphi$. As $f(\pi/2)=f(\pi)=0$,
\eqref{eq:stab-ftc} gives
\begin{align*}
  f(t)&=-\sin t\int_{\pi/2}^t\frac{r_4(u)}{\sin u}\dd u,
    & t&\in[\pi/2,\pi),\\
  f(t)&=-\frac{\cos t}{\cos\varphi}f(T)
    +\sin(T-t)\int_t^{\pi/2}\frac{r_3(u)\dd u}{\sin(T-u)},
    & t&\in[b',\pi/2],\\
  f(t)&=f(b')-\int_t^{b'}f(u+\tfrac\pi2)\dd u+\int_t^{b'}r_2,
    & t&\in[\varphi,b'],\\
  f(t)&=\cos t\Bigl(\frac{f(\varphi)}{\cos\varphi}
    +\int_t^\varphi\frac{r_1(u)}{\cos u}\dd u\Bigr), & t&\in[0,\varphi].
\end{align*}
\ref{prop:stab-coercive:80} Let $m_j=\int\abs{r_j}$ and $M=\sum_jm_j$. On
$[\pi/2,\pi]$, the first formula gives $\abs f\le m_4$, as $0<\sin t\le\sin u$
for $\pi/2\le u\le t$. On $[b',\pi/2]$, the second gives
$\abs{f(t)}\le3\abs{f(T)}+2m_3\le5M$, as $\abs{\cos t}\le\cos\varphi$,
$\sin(T-u)\ge1/2$ and $\abs{f(T)}\le m_4$. On $[\varphi,b']$, the third gives
$\abs f\le5M+2m_4+m_2\le8M$, as $\abs{f(u+\pi/2)}\le m_4$ and
$b'-\varphi\le2$. On $[0,\varphi]$, the fourth gives
$\abs f\le2\abs{f(\varphi)}+2m_1\le18M\le20M$, as
$\cos u\ge\cos\varphi\ge1/2$. The intervals have length at most $2$, so
$m_j^2\le2\int r_j^2$ by the Cauchy--Schwarz inequality, and
$M^2\le4\sum_jm_j^2\le8\sum_j\int r_j^2=16E(f)$.

\ref{prop:stab-coercive:sharp} Let $k(u)=(\sec\varphi+\cos u)/\sin u$. On
$(\pi/2,\pi)$, $r_4=f\cot u-f'$, so $(kf)'=-f-k\,r_4$; integrating this, and
inserting the first two formulas, turns the third into
\[
  f(t)=\int_t^{b'}r_2+\int_{b'}^{\pi/2}\frac{r_3(u)\dd u}{\sin(T-u)}
  +(\sec\varphi-\sin t)\int_{\pi/2}^{\pi/2+t}\frac{r_4(u)\dd u}{\sin u}
  +\int_{\pi/2+t}^{T}k\,r_4 ,
\]
the terms with $f(T)$ cancelling as $k(T)=\tan\varphi$. With $f(\varphi)$ and
$f(T)$ inserted, each formula writes $f(t)$ as a sum of integrals of kernels
against the residuals over disjoint intervals. By the Cauchy--Schwarz
inequality, $f(t)^2$ is at most the sum of the squared $L^2$ norms of the
kernels times $\sum_j\int r_j^2$, and $\int\csc^2=-\cot$, $\int\sec^2=\tan$
and $\int k^2=-(\sec^2\varphi+1)\cot u-2\sec\varphi\csc u-u$ give this sum as
$\mathcal G_\varphi(t)$; on $[\pi/2,\pi)$, for instance, it is
$\sin^2t\,(\cot(\pi/2)-\cot t)=-\sin t\cos t$. Finally
$\mathcal G_\varphi\le2\sec^2\varphi$, as $\tan t\ge0$ on $[0,\varphi]$,
$\cos t\,(2\sec\varphi-\sin t)\le2\sec\varphi\le2\sec^2\varphi$,
$\abs{\sin t\cos t}\le1/2$ and $1/2+2\tan\varphi\le2+2\tan^2\varphi$.

For the last claim, \cref{lem:stab-residuals} gives the hypotheses and
$E(f)=E_\cS(\KG,K)$, which is at most $\abs\Gs-\cQ(\xi)$
(\cref{prop:stab-certificate}).
\end{proof}

\begin{lemma}[downward support values, and distances]\label{lem:stab-lower}
\begin{enumerate}[label=(\alph*)]
\item\label{lem:stab-lower:caps} For $K\in\Kc_{\pi/2}$, $h_K(t)$ is
$h_K(0)\cos t$ if $\sin t\le0\le\cos t$, and $-h_K(\pi)\cos t$ if
$\sin t,\cos t\le0$. So $\dH(K,K')=\max_{[0,\pi]}\abs{h_K-h_{K'}}$ for
$K,K'\in\Kc_{\pi/2}$, and if $f$ is formed with caps of $\Kc_{\pi/2}$ and
$\abs f\le\rho$ on $[0,\pi]$, then $\abs f\le\rho$ on $\R$.
\item\label{lem:stab-lower:close} Convex bodies $K,L$ with $\dH(K,L)\le r$
are $r$-close; and if nonempty compact sets $X,Y$ are $r$-close, then
$\abs{h_X-h_Y}\le r$.
\end{enumerate}
\end{lemma}

\begin{proof}
\ref{lem:stab-lower:caps} Every $p\in K$ has $p_y\ge0$ and
$-h_K(\pi)\le p_x\le h_K(0)$, with equality at $A^-_K(0)=(h_K(0),0)$ and
$C^+_K(\pi/2)=(-h_K(\pi),0)$ (\cref{lem:ends}). This gives $h_K(t)$ for
$\sin t\le0$, where $f(t)$ is then $f(0)\cos t$ or $0$.
\ref{lem:stab-lower:close} For $p\in K$,
$\ip p{u_t}\le h_L(t)+r=h_{L+r\mathbb B(0,1)}(t)$ for all $t$, so
$p\in L+r\mathbb B(0,1)$, and symmetrically. If $p\in X$ has
$\ip p{u_t}=h_X(t)$ and $q\in Y$ has $\norm{p-q}\le r$, then
$h_Y(t)\ge\ip q{u_t}\ge h_X(t)-r$, and symmetrically.
\end{proof}

\begin{proof}[Proof of \cref{thm:stability-cap}]
\ref{thm:stability-cap:wide} Let $f$ be formed with $\KG$ and $K$, so that
$f=h_K-h_{\KG+(s_K,0)}$. By \cref{lem:stab-residuals},
\cref{prop:stab-coercive}\ref{prop:stab-coercive:sharp} applies and gives
$\abs f\le2\sec\varphi\sqrt{E_\cS(\KG,K)}$ on $[0,\pi]$, and
$E_\cS(\KG,K)\le\abs\Gs-\cQ(\xi)$ by \cref{prop:stab-certificate}. By
\cref{lem:stab-lower}\ref{lem:stab-lower:caps} the bound holds on $\R$, so
$\dH(K,\KG+(s_K,0))\le2\sec\varphi\sqrt{\abs\Gs-\cQ(\xi)}$, and
\cref{lem:stab-lower}\ref{lem:stab-lower:close} gives the closeness.
\ref{thm:stability-cap:ki} For $K\in\Ki$, $\xi_K\in\cT\subset\bar\cT$ and
$\cA(K)\le\cQ(\xi_K)$ (\cref{fact:Q}\ref{fact:Q:triple}); so
\ref{thm:stability-cap:wide} applies to $\xi_K$, and
$\abs\Gs-\cQ(\xi_K)\le\abs\Gs-\cA(K)$. \ref{thm:stability-cap:num} As
$\varphi\le0.04$ (\cref{fact:gerver}\ref{fact:gerver:def}),
$\cos\varphi\ge1-\varphi^2/2\ge1249/1250$.
\end{proof}

\subsection{The upper bound near Gerver's cap}\label{sec:stab-local}

Baek proves $\cA(K)\le\cQ(\xi_K)$ for $K\in\Ki$ \baek{Thm.~8.2.4}; this
subsection proves it for the caps of $\Kc_{\pi/2}$ near $\KG$. If
$\dH(K,K')\le\delta$, then $\norm{\xx_K(t)-\xx_{K'}(t)}\le2\delta$, by the
formula for the inner corner (\cref{sec:sofas}). Let $\Gamma$ be as in
\cref{fact:gerver}\ref{fact:gerver:niche}, and let $a_\Gamma=\bD(0)_x$ and
$b_\Gamma=\bB(\pi/2)_x$, the numbers $a<b$ of the proof of
\cref{prop:regular}. The top edge $e_{\KG}(\pi/2)$ of $\KG$ is the segment from
$C^+_{\KG}(0)=\bC(0)=(a_\Gamma,1)$ to $A^-_{\KG}(\pi/2)=\bA(\pi/2)=(b_\Gamma,1)$
(\cref{fact:gerver}\ref{fact:gerver:cap} and the proof of \cref{prop:regular}).
Recall that the width of a cap $K$ is $h_K(0)+h_K(\pi)$; its \emph{cut
half-planes} are $\breve H^{\mathrm R}_K=H^{\mathrm b}_K(\varphi)$ and
$\breve H^{\mathrm L}_K=H^{\mathrm d}_K(\pi/2-\varphi)$ \baek{Def.~8.1.5}.

\begin{proposition}[the upper bound near $\KG$]\label{prop:stab-local}
There is $\delta\in(0,1]$ such that every $K\in\Kc_{\pi/2}$ with
$\dH(K,\KG)\le\delta$ satisfies $\xi_K=(K,B_K,D_K)\in\bar\cT$,
$\cN(K)\subset K$ and $\cA(K)\le\cQ(\xi_K)\le\abs\Gs$, and $K$ and
$\KG+(s_K,0)$ are $80\sqrt{\abs\Gs-\cA(K)}$-close.
\end{proposition}

\begin{lemma}[Gerver's roof]\label{lem:stab-roof}
There are $\bar\gamma<1$, $L\ge0$ and an $L$-Lipschitz
$\gamma:[a_\Gamma,b_\Gamma]\to[0,\bar\gamma]$ with
$\gamma(a_\Gamma)=\gamma(b_\Gamma)=0$ and graph $\Gamma$, such that
$\cN(\KG)=\{(x,y):x\in[a_\Gamma,b_\Gamma],\ 0\le y<\gamma(x)\}$. Moreover
$-h_{\KG}(\pi)<a_\Gamma<b_\Gamma<h_{\KG}(0)$ and
$[a_\Gamma,b_\Gamma]\times[0,1]\subset\KG$.
\end{lemma}

\begin{proof}
The chords of $\Gamma$ have bounded slopes: on $[t_0,t_2]$, $\bD'$ is a
positive multiple of $u_t$ (\cref{fact:gerver}\ref{fact:gerver:niche}) and
$\cos t\ge1/2$; likewise for $\bB$ on $[t_3,t_5]$; and on $[t_1,t_4]$ the
abscissa of $\xx$ has the derivative $\alpha\cos t-\beta\sin t$, which is
negative (\cref{lem:gerver-width}), hence bounded away from $0$. As the three
pieces occupy consecutive intervals of abscissas and meet at their ends,
$\Gamma$ is the graph of an $L$-Lipschitz $\gamma$ on $[a_\Gamma,b_\Gamma]$,
with $\gamma(a_\Gamma)=\bD(0)_y=0$ and $\gamma(b_\Gamma)=\bB(\pi/2)_y=0$. The
proof of \cref{prop:regular} gives the heights in $[0,1)$, with a largest,
$\bar\gamma$, and $[a_\Gamma,b_\Gamma]\times[0,1]\subset\KG$; the niche is that
of \cref{fact:gerver}\ref{fact:gerver:niche}. As $\sigma_{\KG}$ has no atom at
$\pi$ (\cref{fact:arms}\ref{fact:arms:regular}), the edge $e_{\KG}(\pi)$ is the
point $(-h_{\KG}(\pi),0)$ (\cref{lem:ends}). The point $(a_\Gamma,1)$ of $\KG$
is not this point, so $\ip{(a_\Gamma,1)}{u_\pi}<h_{\KG}(\pi)$, that is,
$-h_{\KG}(\pi)<a_\Gamma$. Likewise $b_\Gamma<h_{\KG}(0)$.
\end{proof}

\begin{lemma}[persistence on edges]\label{lem:stab-faces}
Let $K_0\in\Kc_{\pi/2}$ and $\mathbb B=\mathbb B(0,1)$.
\begin{enumerate}[label=(\alph*)]
\item\label{lem:stab-faces:fixed} If $I'\subset[0,\pi]$ is compact and $F$ is
continuous on $(K_0+\mathbb B)\times I'$, with $F(p,t)>0$ for $t\in I'$ and
$p\in e_{K_0}(t)$, then there is $\delta\in(0,1]$ such that $F(p,t)>0$ for all
$K\in\Kc_{\pi/2}$ with $\dH(K,K_0)\le\delta$, $t\in I'$ and $p\in e_K(t)$.
\item\label{lem:stab-faces:angle} If $F$ is continuous on $\R^2$ and $F>0$ on
$e_{K_0}(t_0)$, then there are $\delta,\rho>0$ such that $F(p)>0$ for all such
$K$, $t\in[0,\pi]$ with $\abs{t-t_0}\le\rho$, and $p\in e_K(t)$.
\item\label{lem:stab-faces:arms} If $K_0\in\Ki$ and $0<s_1\le s_2<\pi/2$,
there are $c,\delta>0$ such that every $K\in\Kc_{\pi/2}$ with
$\dH(K,K_0)\le\delta$ has, for $t\in[s_1,s_2]$,
$h_K(t+\pi/2)-\ip p{v_t}\ge1+c$ for $p\in e_K(t)$ and
$h_K(t)-\ip p{u_t}\ge1+c$ for $p\in e_K(t+\pi/2)$. So $f^\pm_K,g^\pm_K\ge1+c$
on $[s_1,s_2]$, and there the abscissa of $\xx_K$ has right derivatives at most
$-c$.
\end{enumerate}
\end{lemma}

\begin{proof}
\ref{lem:stab-faces:fixed} The function
$\Theta(p,t)=\operatorname{dist}(p,K_0)+\abs{\ip p{u_t}-h_{K_0}(t)}$ vanishes
exactly when $p\in e_{K_0}(t)$, so it has a positive minimum $m$ on the compact
set where $F\le0$; let $\delta=\min(1,m/4)$. A cap $K$ with
$\dH(K,K_0)\le\delta$ lies in $K_0+\mathbb B$ and is $\delta$-close to $K_0$
(\cref{lem:stab-lower}), so $\Theta(p,t)\le2\delta<m$ for $p\in e_K(t)$.
\ref{lem:stab-faces:angle} The proof of \ref{lem:stab-faces:fixed} applies to
$\Theta(p)=\operatorname{dist}(p,K_0)+\abs{\ip p{u_{t_0}}-h_{K_0}(t_0)}$, which
is at most $2\delta+4R\abs{t-t_0}$ for $p\in e_K(t)$, where $R$ bounds the
norms on $K_0+\mathbb B$. \ref{lem:stab-faces:arms} By condition~(3) of the
injectivity condition and \cref{fact:arms}\ref{fact:arms:regular}, $f_{K_0}$
and $g_{K_0}$ are continuous and exceed $1$ on $(0,\pi/2)$, hence exceed $1+m$
on $[s_1,s_2]$ for some $m>0$, and $e_{K_0}(t)$, $e_{K_0}(t+\pi/2)$ are the
points $A_{K_0}(t)$, $C_{K_0}(t)$. Part \ref{lem:stab-faces:fixed} for
$F(p,t)=h_{K_0}(t+\pi/2)-\ip p{v_t}-1-m/2$ and its analogue, with
$\delta\le m/4$, gives $c=m/4$. The right derivative of $\xx_K$ is
$-(f^+_K-1)u_t+(g^+_K-1)v_t$ (\cref{fact:arms}\ref{fact:arms:deriv}), with
abscissa at most $-c(\cos t+\sin t)\le-c$.
\end{proof}

\begin{lemma}[the niche of a nearby cap]\label{lem:stab-niche}
\begin{enumerate}[label=(\alph*)]
\item\label{lem:stab-niche:x} For every $\eta>0$ there is $\delta\in(0,1]$
such that every $K\in\Kc_{\pi/2}$ with $\dH(K,\KG)\le\delta$ has
$a_\Gamma-\eta\le p_x\le b_\Gamma+\eta$ for all $p\in\cN(K)$.
\item\label{lem:stab-niche:sub} There is $\delta\in(0,1]$ such that
$\cN(K)\subset K$ for every $K\in\Kc_{\pi/2}$ with $\dH(K,\KG)\le\delta$.
\item\label{lem:stab-niche:width} Let $K\in\Kc_{\pi/2}$. If
$\dH(K,\KG)\le1/20$, then $K$ has width at least $21/10$. If $K$ has width at
least $21/10$, then for every $\varphi\in[0.039,0.04]$, Gerver's angle among
them, $W_K(\varphi)$ and $Z_K(\pi/2-\varphi)$ lie on the bottom edge of $K$,
strictly between its ends, and
$K\cap H^{\mathrm b}_K(\varphi)\cap H^{\mathrm d}_K(\pi/2-\varphi)=\emptyset$.
\end{enumerate}
\end{lemma}

\begin{proof}
\ref{lem:stab-niche:x} A point $p\in\cN(K)$ has $p_y\ge0$ and lies in
$Q^-_K(t)$ for some $t\in(0,\pi/2)$, so $Z_K(t)_x<p_x<W_K(t)_x$, where
$Z_K(t)_x=(1-h_K(t+\pi/2))/\sin t$ and $W_K(t)_x=(h_K(t)-1)/\cos t$. For
$\KG$, $\xx(t)_y=\sin t\cos t\,(W_{\KG}(t)_x-Z_{\KG}(t)_x)$ is positive (by
\cref{fact:height} on $[t_1,t_4]$, and elsewhere as $\xx(t)$ lies above
$\bD(t)$ or $\bB(t)$), so the open segment of the $x$-axis between
$Z_{\KG}(t)$ and $W_{\KG}(t)$ lies in $\cN(\KG)$, and
$a_\Gamma\le Z_{\KG}(t)_x$ and $W_{\KG}(t)_x\le b_\Gamma$
(\cref{lem:stab-roof}). As the top edge of $\KG$ is
$[(a_\Gamma,1),(b_\Gamma,1)]$, \cref{lem:stab-faces}\ref{lem:stab-faces:angle}
at $t_0=\pi/2$ gives $\rho\in(0,\pi/4]$ and $\delta_0$ such that, if
$\delta\le\delta_0$ and $t\le\rho$, each $q\in e_K(t+\pi/2)$ has
$q_x>a_\Gamma-\eta$, and $Z_K(t)_x=q_x+(1-q_y\cos t)/\sin t\ge q_x$. For
$t>\rho$, $\abs{Z_K(t)_x-Z_{\KG}(t)_x}\le\delta/\sin\rho\le\eta$ if
$\delta\le\eta\sin\rho$. So $p_x>a_\Gamma-\eta$, and symmetrically
$p_x<b_\Gamma+\eta$.

\ref{lem:stab-niche:sub} Let $p\in\cN(K)$, so that $p\in Q^-_K(t)$ for some
$t\in(0,\pi/2)$. Let $M'=\max\xx_y<1$ (\cref{fact:height}) and
$\delta\le(1-M')/4$. As $(0,1)=u_t\sin t+v_t\cos t$,
$p_y<\xx_K(t)_y\le\xx(t)_y+2\delta\le y'=(1+M')/2$. By compactness, there is
$m>0$ with $h_{\KG}(t)-\ip q{u_t}\ge m$ for all
$q\in[a_\Gamma,b_\Gamma]\times[0,y']$ and $t\in[0,\pi]$, as $(q_x,1)\in\KG$
and $-h_{\KG}(\pi)<a_\Gamma<b_\Gamma<h_{\KG}(0)$. With $\delta\le m/4$ and
\ref{lem:stab-niche:x} for $\eta=m/4$, $p$ lies within $m/4$ of such a $q$, so
$\ip p{u_t}\le h_{\KG}(t)-3m/4\le h_K(t)$ for $t\in[0,\pi]$, and $p\in K$ by
\eqref{eq:cap-fan}.

\ref{lem:stab-niche:width} By \cref{lem:stab-lower}, a cap
$K'\in\Kc_{\pi/2}$ lies in $[-h_{K'}(\pi),h_{K'}(0)]\times[0,1]$, so its area
is at most its width. Hence $\KG$ has width at least $\abs{\KG}\ge\abs\Gs>11/5$
(\cref{fact:gerver}\ref{fact:gerver:def}), and $K$, whose support values
differ from those of $\KG$ by at most $1/20$, has width at least
$11/5-2/20=21/10$. As $h_K(\varphi)\ge h_K(0)\cos\varphi$ and
$(21/10)\cos\varphi>1$, $W_K(\varphi)_x>-h_K(\pi)$; and
$W_K(\varphi)_x<h_K(0)$, as
$h_K(\varphi)\le h_K(0)\cos\varphi+h_K(\pi/2)\sin\varphi<h_K(0)\cos\varphi+1$
by the sublinearity of $h_K$. Likewise for $Z_K(\pi/2-\varphi)$. A point $p$
of $K$ in $H^{\mathrm b}_K(\varphi)\cap H^{\mathrm d}_K(\pi/2-\varphi)$ would
satisfy
\[
  2\sin\varphi\ge\ip p{u_\varphi+u_{\pi-\varphi}}
  \ge h_K(\varphi)+h_K(\pi-\varphi)-2\ge(21/10)\cos\varphi-2,
\]
where the first inequality holds as $u_\varphi+u_{\pi-\varphi}=(0,2\sin\varphi)$
and $p_y\le1$, the second by the definition of the two half-planes, and the
third as $h_K(\varphi)\ge h_K(0)\cos\varphi$ and
$h_K(\pi-\varphi)\ge h_K(\pi)\cos\varphi$; as $\varphi\le0.04$, this is false.
\end{proof}

\begin{lemma}[the canonical triple and the cuts]\label{lem:stab-triple}
There is $\delta\in(0,1]$ such that every $K\in\Kc_{\pi/2}$ with
$\dH(K,\KG)\le\delta$ satisfies:
\begin{enumerate}[label=(\alph*)]
\item\label{lem:stab-triple:wide} $\xi_K\in\bar\cT$;
\item\label{lem:stab-triple:cut} $\xx_K(t)\notin\breve H^{\mathrm R}_K$ for
$t\in(\varphi,\pi/2]$, and $\xx_K(t)\notin\breve H^{\mathrm L}_K$ for
$t\in[0,\pi/2-\varphi)$;
\item\label{lem:stab-triple:height} $\xx_K(t)_y>0$ for
$t\in[\varphi,\pi/2-\varphi]$.
\end{enumerate}
\end{lemma}

\begin{proof}
Take $c$ and $\delta\le1/20$ from
\cref{lem:stab-faces}\ref{lem:stab-faces:arms} for $K_0=\KG$ and
$[\varphi,\pi/2-\varphi]$; so $g^+_K(\varphi)>1$, $f^-_K(\pi/2-\varphi)>1$,
and $K$ has width at least $21/10$
(\cref{lem:stab-niche}\ref{lem:stab-niche:width}).
\ref{lem:stab-triple:wide} The inequalities defining $\bar\cT$ hold for $B_K$
and $D_K$ by definition. For the equalities at the ends $\pi/2$ and $0$, $B_K$
contains $(h_K(0),0)$ and $D_K$ contains $(-h_K(\pi),0)$: for
$t\in[\varphi,\pi/2]$, $h_K(t)\le h_K(0)\cos t+h_K(\pi/2)\sin t\le
h_K(0)\cos t+1$ by the sublinearity of $h_K$, and likewise for $D_K$. For the
equality at $\varphi$ we follow Baek's proof of \baek{Lemma~8.1.7}:
$b_K(\varphi)$ contains $W_K(\varphi)\in K$ (\cref{lem:stab-niche}), and the
point $p$ of $K\cap b_K(\varphi)$ farthest in the direction $v_\varphi$ lies in
some $e_K(\vartheta)$, $\vartheta\in[\varphi,\varphi+\pi]$. If
$\vartheta>\varphi+\pi/2$, then
$\ip p{u_\varphi}\le\ip{v^+_K(\varphi+\pi/2)}{u_\varphi}$, that is,
$g^+_K(\varphi)\le1$. So $\vartheta\le\varphi+\pi/2$, and the sublinearity of
$h_K$ between $\varphi$, $\vartheta$ and $\pi/2$ (where $h_K=1\ge p_y$) gives
$\ip p{u_t}\ge h_K(t)-1$ for $t\in[\varphi,\pi/2]$, so $p\in B_K$. The mirror
argument with $f^-_K(\pi/2-\varphi)>1$ treats $D_K$.
\ref{lem:stab-triple:cut} Let $t_*=(\varphi+\pi/2)/2$. For
$t\in(\varphi,t_*]$, the right derivative of $\ip{\xx_K}{u_\varphi}$ is
$-(f^+_K-1)\cos(t-\varphi)-(g^+_K-1)\sin(t-\varphi)\le-c$, so
$\ip{\xx_K(t)}{u_\varphi}<\ip{\xx_K(\varphi)}{u_\varphi}=h_K(\varphi)-1$. On
$[t_*,\pi/2]$, $\KG$ satisfies the strict inequality
$\ip{\xx(t)}{u_\varphi}<h_{\KG}(\varphi)-1$ \baek{Lemma~8.1.6} with a positive
minimal gap, and this persists for small $\delta$, as
$\ip{\xx_K(t)}{u_\varphi}$ and $h_K(\varphi)$ differ from their values for
$\KG$ by at most $2\delta$ and $\delta$. The left cut is symmetric.
\ref{lem:stab-triple:height} $\xx(t)_y$ has a positive minimum $m'$ on
$[\varphi,\pi/2-\varphi]$ (\cref{fact:height}), and
$\xx_K(t)_y\ge\xx(t)_y-2\delta>0$ if $\delta<m'/2$.
\end{proof}

\begin{lemma}[the three regions of the niche]\label{lem:stab-areabound}
Let $\varphi$ be any number in $[0.039,0.04]$, and form $\bar\cT$, $\xi_K$,
$I$, $\cQ$ and the cut half-planes with it. Let $K\in\Kc_{\pi/2}$ have width
at least $21/10$, $\xi_K\in\bar\cT$, $\cN(K)\subset K$, the properties
\ref{lem:stab-triple:cut} and \ref{lem:stab-triple:height} of
\cref{lem:stab-triple}, and those of
\cref{lem:stab-faces}\ref{lem:stab-faces:arms} on $[\varphi,\pi/2-\varphi]$
for some $c>0$. Then $\cA(K)\le\cQ(\xi_K)$.
\end{lemma}

\begin{proof}
We follow Baek's proof, which bounds the areas of three regions of the niche
\baek{Lemma~8.2.2, Lemma~8.2.3, Thm.~8.2.4}. Write $B=B_K$, $D=D_K$,
$x^{\mathrm R}=\xx_K(\varphi)$, $x^{\mathrm L}=\xx_K(\pi/2-\varphi)$,
$W=W_K(\varphi)$ and $Z=Z_K(\pi/2-\varphi)$. As $\cN(K)\subset K$,
\cref{lem:stab-niche}\ref{lem:stab-niche:width} splits $\cN(K)$ into its parts
in $\breve H^{\mathrm R}_K$, in $\breve H^{\mathrm L}_K$ and outside both. For
$t\in(\varphi,\pi/2]$,
$\breve H^{\mathrm R}_K\cap Q^-_K(t)=\breve H^{\mathrm R}_K\setminus
H^{\mathrm b}_K(t)$, as $\xx_K(t)\notin\breve H^{\mathrm R}_K$ and
$u_\varphi=\cos(t-\varphi)u_t-\sin(t-\varphi)v_t$. So the points of $K$
outside $B$ in the supporting half-planes of $B$ with normals $\pi+\varphi$
and $3\pi/2$ lie in $\cN(K)\cap\breve H^{\mathrm R}_K$; among them are the
points of the region between the arc of $B$ with normals in
$(\pi+\varphi,3\pi/2)$ and its end tangents, which meet at $W\in K$. This
region has area
$\mathcal J([X_B,W])-\frac12\int_{(\pi+\varphi,3\pi/2)}h_B\dd\sigma_B$
\baek{Lemma~7.3.5}, so $\abs{\cN(K)\cap\breve H^{\mathrm R}_K}$ is at least
this area. Likewise
$\abs{\cN(K)\cap\breve H^{\mathrm L}_K}\ge\mathcal J([Z,Y_D])
-\frac12\int_{(3\pi/2,3\pi/2+\varphi^{\mathrm L})}h_D\dd\sigma_D$. By
\cref{lem:stab-faces}\ref{lem:stab-faces:arms}, the abscissa of $\xx_K$
decreases on $I$ with right derivatives at most $-c$, so $\xx_K|_I$ is the
graph of a Lipschitz function of the abscissa. The region $\Omega$ between this
graph and the $x$-axis and the triangles $T^{\mathrm R}$ with vertices
$(x^{\mathrm R}_x,0)$, $W$, $x^{\mathrm R}$ and $T^{\mathrm L}$ with vertices
$Z$, $(x^{\mathrm L}_x,0)$, $x^{\mathrm L}$ are disjoint, and lie in $\cN(K)$
outside the cut half-planes: points below $\xx_K(t)$ lie in $Q^-_K(t)$, and
\cref{lem:stab-triple}\ref{lem:stab-triple:cut} applies. By changes of
variables and integration by parts with right derivatives,
\begin{align*}
  \abs{\Omega\cup T^{\mathrm R}\cup T^{\mathrm L}}
  &=-\int_I(\xx_K)'_x(\xx_K)_y\dd t+\tfrac12\tan\varphi\,
    \bigl((x^{\mathrm R}_y)^2+(x^{\mathrm L}_y)^2\bigr)\\
  &=\mathcal J([W,x^{\mathrm R}])+\mathcal J(\xx_K|_I)
    +\mathcal J([x^{\mathrm L},Z]).
\end{align*}
Adding the three bounds and joining the collinear points
$x^{\mathrm R},X_B,W$ on $b_K(\varphi)$ and $Y_D,x^{\mathrm L},Z$ on
$d_K(\pi/2-\varphi)$ gives $\abs{\cN(K)}\ge\abs K-\cQ(\xi_K)$.
\end{proof}

\begin{proof}[Proof of \cref{prop:stab-local}]
Let $\delta$ be the least of $1/20$ and of those of \cref{lem:stab-triple},
\cref{lem:stab-niche}\ref{lem:stab-niche:sub} and
\cref{lem:stab-faces}\ref{lem:stab-faces:arms} for $K_0=\KG$ and
$[\varphi,\pi/2-\varphi]$; recall $\KG\in\Ki$
(\cref{fact:gerver}\ref{fact:gerver:cap}) and $\varphi\in[0.039,0.04]$
(\cref{fact:gerver}\ref{fact:gerver:def}). By
\cref{lem:stab-niche}\ref{lem:stab-niche:width}, $K$ has width at least
$21/10$. Then
\cref{lem:stab-triple,lem:stab-niche,lem:stab-areabound,lem:stab-variation}
give $\xi_K\in\bar\cT$, $\cN(K)\subset K$ and
$\cA(K)\le\cQ(\xi_K)\le\abs\Gs$, and \cref{prop:stab-coercive,lem:stab-lower}
show that $K$ and $\KG+(s_K,0)$ are $80\sqrt{\abs\Gs-\cQ(\xi_K)}$-close, hence
$80\sqrt{\abs\Gs-\cA(K)}$-close, as $\cA(K)\le\cQ(\xi_K)$.
\end{proof}

\subsection{The cost of a missing final angle}\label{sec:stab-angle}

For $K\in\Kc_{\pi/2}$ and $t\in\R$, the coordinates of $p-\xx_K(t)$ in the
frame $(u_t,v_t)$ are $\ell^{\mathrm b}_K(t,p)=\ip p{u_t}-h_K(t)+1$ and
$\ell^{\mathrm d}_K(t,p)=\ip p{v_t}-h_K(t+\pi/2)+1$; we write $\ell_K(t,p)$
for the pair $(\ell^{\mathrm b}_K(t,p),\ell^{\mathrm d}_K(t,p))$. Then
$Q^-_K(t)$ is the set where both coordinates are negative. For
$\omega\in[0,\pi/2]$, let $\cN_\omega(K)$ be the set of $p$ with $p_y\ge0$ and
$p\in Q^-_K(t)$ for some $t\in(0,\omega)$, so that $\cN_{\pi/2}(K)=\cN(K)$,
and let $O_\omega(K)=(K\setminus\cN_\omega(K))\setminus(K\setminus\cN(K))$ be
the \emph{omitted wedges}. A set $S$ \emph{satisfies the constraints of
$(K,\omega)$} if $S\subset K$,
$\max(\ell^{\mathrm b}_K(t,p),\ell^{\mathrm d}_K(t,p))\ge0$ for $p\in S$ and
$t\in[0,\omega]$, and $\ell^{\mathrm b}_K(\omega,p)\ge0$ for $p\in S$.

\begin{proposition}[the terminal comparison]\label{prop:stab-terminal}
There are $c>0$, $\delta\in(0,1]$, $\psi_0>0$ and $R\ge1$ such that for every
$K\in\Kc_{\pi/2}$ with $\dH(K,\KG)\le\delta$, every $\omega\in[0,\pi/2]$ with
$\psi=\pi/2-\omega\le\psi_0$ and every measurable set $S$ that satisfies the
constraints of $(K,\omega)$, with $\eps=\abs\Gs-\abs S$ and
$U=K\setminus\cN(K)$:
\begin{enumerate}[label=(\alph*)]
\item\label{prop:stab-terminal:area} $\abs S\le\cA(K)-c\psi$,
$\psi\le\eps/c$ and $\abs\Gs-\cA(K)\le\eps$;
\item\label{prop:stab-terminal:diff} $\abs{S\setminus U}\le\eps$ and
$\abs{U\setminus S}\le2\eps$;
\item\label{prop:stab-terminal:hall}
$\max(\ell^{\mathrm b}_K(t,p),\ell^{\mathrm d}_K(t,p))\ge-4R\psi$ for $p\in S$
and $t\in(0,\pi/2)$.
\end{enumerate}
\end{proposition}

\begin{lemma}[a rectangle missed by the sofa]\label{lem:stab-floor}
There are $c_0,\delta,\psi_0>0$ such that for every $K\in\Kc_{\pi/2}$ with
$\dH(K,\KG)\le\delta$, every $\psi\in(0,\psi_0]$ and every set $S$ with
$\ell^{\mathrm b}_K(\pi/2-\psi,p)\ge0$ for $p\in S$, there is a measurable
$F\subset K\setminus\cN(K)$ with $F\cap S=\emptyset$ and $\abs F=c_0\psi$.
\end{lemma}

\begin{proof}
Let $x_0=-h_{\KG}(\pi)$, $d=(a_\Gamma-x_0)/2>0$ (\cref{lem:stab-roof}),
$l=x_0+d/2$ and $r=a_\Gamma-d$. As $(x_0,0),(a_\Gamma,1)\in\KG$, the rectangle
$[l,r]\times[0,1/4]$ lies in $\KG$, so by compactness there is $m>0$ such that
the points of $[l,r]\times[0,1/8]$ have distance at least $m$ from the
supporting lines of $\KG$ with normals in $[0,\pi]$. For $\delta\le m/2$ this
rectangle lies in $K$ by \eqref{eq:cap-fan}, and by
\cref{lem:stab-niche}\ref{lem:stab-niche:x} with $\eta=d/4$ it misses
$\cN(K)$. By \cref{lem:stab-faces}\ref{lem:stab-faces:angle} for the top edge
of $\KG$, $q=v^+_K(\pi/2)$ has $q_y=1$ and $q_x>a_\Gamma-d/4$, for small
$\delta$. Let $c_0=(r-l)d/8$, $\psi_0=\min(1,d/4,1/d)$ and
$F=[l,r]\times[0,\psi d/8]$. For $p\in F$, as
$h_K(\pi/2-\psi)\ge q_x\sin\psi+\cos\psi$, $q_x-p_x\ge3d/4$,
$\sin\psi\ge\psi/2$ and $1-\cos\psi\le\psi^2/2\le\psi d/8$,
\[
  \ell^{\mathrm b}_K(\tfrac\pi2-\psi,p)\le-(q_x-p_x)\sin\psi+p_y\cos\psi
    +1-\cos\psi\le-\tfrac{3\psi d}8+\tfrac{\psi d}8+\tfrac{\psi d}8<0,
\]
so $p\notin S$.
\end{proof}

\begin{lemma}[the omitted wedges]\label{lem:stab-wedges}
For all $\eta>0$ and $R\ge1$ there are $\delta\in(0,1]$ and $\psi_0>0$ such
that $\abs{O_\omega(K)}\le4\eta(3R+1)(\pi/2-\omega)$ for every
$K\in\Kc_{\pi/2}$ with $\dH(K,\KG)\le\delta$ and $\norm p\le R$ on $K$, and
every $\omega\in[0,\pi/2]$ with $\pi/2-\omega\le\psi_0$.
\end{lemma}

\begin{proof}
A point $p\in O_\omega(K)$ has $p_y\ge0$ and both coordinates of
$\ell_K(t,p)$ negative at some $t\in[\omega,\pi/2)$, so $p_y<\xx_K(t)_y$
(proof of \cref{lem:stab-niche}). As $h_K$ is $2R$-Lipschitz, $h_K(\pi/2)=1$,
$\abs{h_K}\le R$ and $\cos t\le\pi/2-t$,
$\xx_K(t)_y=(h_K(t)-1)\sin t+(h_K(t+\pi/2)-1)\cos t\le(3R+1)(\pi/2-\omega)$.
By \cref{lem:stab-niche}\ref{lem:stab-niche:x},
$a_\Gamma-\eta\le p_x\le b_\Gamma+\eta$. Every $x\in(a_\Gamma,b_\Gamma)$ is the
abscissa of a point of $\Gamma$ other than $\bD(t_0)$ and $\bB(t_5)$, and these
points have positive height: $\bD_y$ increases strictly from $0$ on
$[t_0,t_2]$, $\bB_y$ decreases strictly to $0$ on $[t_3,t_5]$
(\cref{fact:gerver}\ref{fact:gerver:niche}), and $\xx_y>0$ on $[t_1,t_4]$
(\cref{fact:height}). So the segment $(a_\Gamma,b_\Gamma)\times\{0\}$ lies in
$\cN(\KG)$, and by compactness there is $\nu>0$ such that each
$x\in[a_\Gamma+\eta,b_\Gamma-\eta]$ has a $t\in(\nu,\pi/2-\nu)$ at which both
coordinates of $\ell_{\KG}(t,(x,0))$ are less than $-\nu$. If $\delta\le\nu/4$,
the points of $[a_\Gamma+\eta,b_\Gamma-\eta]\times[0,\delta]$ then lie in
$\cN_\omega(K)$ when $\omega>\pi/2-\nu$. With
$\psi_0=\min(\nu/2,\delta/(3R+1))$, $O_\omega(K)$ lies in
$([a_\Gamma-\eta,a_\Gamma+\eta]\cup[b_\Gamma-\eta,b_\Gamma+\eta])
\times[0,(3R+1)(\pi/2-\omega)]$.
\end{proof}

\begin{proof}[Proof of \cref{prop:stab-terminal}]
Caps with $\dH(K,\KG)\le1$ lie in $\KG+\mathbb B(0,1)\subset\mathbb B(0,R)$
for some $R\ge1$. Let $c_0$ be as in \cref{lem:stab-floor}, $c=c_0/2$ and
$\eta=c_0/(16(3R+1))$, and let $\delta,\psi_0$ be the least of those of
\cref{lem:stab-floor,lem:stab-wedges,prop:stab-local}. By
\cref{prop:stab-local}, $\abs U=\cA(K)\le\abs\Gs$. The set
$V=K\setminus\cN_\omega(K)$ contains $U$ and, by the constraints, $S$, and
$\abs{V\setminus U}=\abs{O_\omega(K)}\le c\psi$. \Cref{lem:stab-floor}, or
$F=\emptyset$ if $\psi=0$, gives $F\subset U$, disjoint from $S$, with
$\abs F=2c\psi$, so
\[
  \abs S+2c\psi\le\abs V=\abs U+\abs{V\setminus U}\le\cA(K)+c\psi ,
\]
where the first inequality holds as $S$ and $F$ are disjoint subsets of $V$,
and the last as $\abs U=\cA(K)$ and $\abs{V\setminus U}\le c\psi$. With
$\cA(K)\le\abs\Gs$ this gives \ref{prop:stab-terminal:area}, then
$\abs{S\setminus U}\le\abs{V\setminus U}\le c\psi\le\eps$ and
$\abs{U\setminus S}=\abs U-\abs S+\abs{S\setminus U}\le2\eps$. For
\ref{prop:stab-terminal:hall}, $\ell^{\mathrm b}_K(\cdot,p)$ and
$\ell^{\mathrm d}_K(\cdot,p)$ are $4R$-Lipschitz for $p\in K$, and their
maximum is nonnegative at $\omega$.
\end{proof}

\subsection{From the cap back to the sofa}\label{sec:stab-recovery}

\begin{lemma}[a right-angle cap of a sofa]\label{lem:stab-sofacap}
Let $S$ be a moving sofa with rotation angle $\omega$ and $h_S(\pi/2)=1$, and
let $K=\{p:p_y\ge0,\ \ip p{u_t}\le h_S(t)\text{ for }t\in[0,\pi]\}$.
\begin{enumerate}[label=(\alph*)]
\item\label{lem:stab-sofacap:sofa} $S\subset\R\times[0,1]$,
$h_S(\omega)+h_S(\omega+\pi)\le1$, and $S\subset L_S(t)$ for
$t\in[0,\omega]$.
\item\label{lem:stab-sofacap:cap} $K$ is a cap in $\Kc_{\pi/2}$ with
$S\subset K$ and $h_K=h_S$ on $[0,\pi]$, and, if $\omega\in[0,\pi/2]$, $S$
satisfies the constraints of $(K,\omega)$. If $S$ and $\Gs$ are
$\delta$-close, then $\dH(K,\KG)\le\delta$.
\end{enumerate}
\end{lemma}

\begin{proof}
\ref{lem:stab-sofacap:sofa} As in the proof of \cref{cor:translate},
$p_y-q_y\le1$ for $p,q\in S$, so $S\subset\R\times[0,1]$, and the end position
$\Phi_1(S)\subset V_L$ gives $\ip{q-p}{u_\omega}\le1$. For $t\in[0,\omega]$,
some $s$ has $\theta(s)=-t$, and $\Phi_s(S)\subset L$ puts $S$ in a translate
of $R_tL$, hence in $L_S(t)$ \baek{Prop.~2.2.3}.
\ref{lem:stab-sofacap:cap} $K$ is closed, convex, contained in
$[-h_S(\pi),h_S(0)]\times[0,1]$, contains $S$ and, with each point $p$, the
point $(p_x,0)$; as an intersection of half-planes with normals in
$[0,\pi]\cup\{3\pi/2\}$, it is a cap, with $h_K=h_S$ on $[0,\pi]$. If
$\omega\le\pi/2$ and $t\in[0,\omega]$, then $t$ and $t+\pi/2$ lie in
$[0,\pi]$, so $\xx_K(t)=\xx_S(t)$; as $S\subset L_S(t)$, at each point of $S$
one of $\ell^{\mathrm b}_K(t,\cdot)$ and $\ell^{\mathrm d}_K(t,\cdot)$ is
nonnegative. Also $\ell^{\mathrm b}_K(\omega,\cdot)\ge0$ on $S$ by
\ref{lem:stab-sofacap:sofa}. Finally $h_{\KG}=h_\Gs$ on $[0,\pi]$
(\cref{fact:monotone}\ref{fact:monotone:cap},
\cref{fact:gerver}\ref{fact:gerver:cap}), and \cref{lem:stab-lower} gives
$\dH(K,\KG)=\max_{[0,\pi]}\abs{h_S-h_\Gs}\le\delta$.
\end{proof}

\begin{lemma}[margins of Gerver's sofa]\label{lem:stab-margins}
\begin{enumerate}[label=(\alph*)]
\item\label{lem:stab-margins:outer} There is $d_0>0$ such that
$h_{\KG}(t)-\ip p{u_t}\ge d_0$ for $p\in\cN(\KG)$ and $t\in[0,\pi]$.
\item\label{lem:stab-margins:roof} There are $c_1,\tau>0$ such that every
$p\in\cN(\KG)$ has a $t\in(0,\pi/2)$ at which both coordinates of
$\ell_{\KG}(t,p)$ are at most $-\min(c_1(\gamma(p_x)-p_y),\tau)$.
\item\label{lem:stab-margins:balls} There are $\kappa,r_0>0$ such that for
every $p\in\Gs$ and $\rho\in(0,r_0]$ some disk $\mathbb B(z,\kappa\rho)$ lies
in $\Gs\cap\mathbb B(p,\rho)$.
\end{enumerate}
\end{lemma}

\begin{proof}
\ref{lem:stab-margins:outer} By \cref{lem:stab-roof},
$\cN(\KG)\subset[a_\Gamma,b_\Gamma]\times[0,\bar\gamma]$, and
$h_{\KG}(t)-\ip p{u_t}>0$ on this rectangle for $t\in[0,\pi]$: at $0$ and
$\pi$ as $-h_{\KG}(\pi)<a_\Gamma<b_\Gamma<h_{\KG}(0)$, and otherwise as
$(p_x,1)\in\KG$ makes it at least $(1-p_y)\sin t$; compactness gives $d_0$.
\ref{lem:stab-margins:roof} Let $q=(p_x,\gamma(p_x))\in\Gamma$ and
$d=q_y-p_y$. Moving down by $d$ lowers $\ell^{\mathrm b}_{\KG}(t,\cdot)$ by
$d\sin t$ and $\ell^{\mathrm d}_{\KG}(t,\cdot)$ by $d\cos t$. If $q=\xx(t)$,
$t\in[t_1,t_4]$, both coordinates of $\ell_{\KG}(t,q)$ vanish, and
$\min(\sin t,\cos t)$ is bounded below. For $q=\bB(t)$, $t\in[t_3,t_5)$,
$\ell^{\mathrm b}_{\KG}(t,q)=0$, $\sin t\ge\sin t_3$, and
$\ell^{\mathrm d}_{\KG}(t,q)=\alpha(t)$ is at most a negative constant, as
$\alpha<0$ on $(0,\pi/2)$ and
$\bB(\pi/2)_x=\xx(\pi/2)_x-\alpha(\pi/2)>\xx(t_1)_x>\xx(\pi/2)_x$. The case
$q=\bD(t)$ is symmetric, with $\beta$. \ref{lem:stab-margins:balls} By
\cref{lem:stab-roof}, $\Gs$ is the union of the convex bodies
$\KG\cap\{x\le a_\Gamma\}$, $\KG\cap\{x\ge b_\Gamma\}$ and
$[a_\Gamma,b_\Gamma]\times[\bar\gamma,1]$, which have interior points, and of
$\{(x,y):x\in[a_\Gamma,b_\Gamma],\ \gamma(x)\le y\le(1+\bar\gamma)/2\}$. A
convex body $C$ with $\mathbb B(z_0,s)\subset C\subset\mathbb B(z_0,D)$ has the
property with $\kappa=s/(D+s)$ and $r_0=D+s$: the homothety of ratio
$\rho/(D+s)$ about $p$ maps $\mathbb B(z_0,s)$ into $C\cap\mathbb B(p,\rho)$.
For a point $p$ of the last set, $\rho\le\min(b_\Gamma-a_\Gamma,
(1-\bar\gamma)/2)$ and $w=\rho/(4(L+2))$, the disk $\mathbb B(z,w)$,
$z=(p_x\pm2w,p_y+(3L+2)w)$ with the sign towards the farther end of
$[a_\Gamma,b_\Gamma]$, lies in $\Gs\cap\mathbb B(p,\rho)$, as $\gamma$ is
$L$-Lipschitz.
\end{proof}

\begin{lemma}[the two directed distances]\label{lem:stab-directed}
Let $K_0,K\in\Kc_{\pi/2}$ and $\delta\ge0$ with $\dH(K,K_0)\le\delta$, and let
$U_0=K_0\setminus\cN(K_0)$ and $U=K\setminus\cN(K)$.
\begin{enumerate}[label=(\alph*)]
\item\label{lem:stab-directed:to} Suppose that $K_0$ has, in place of $\KG$,
the properties of \cref{lem:stab-roof} (with some
$a_\Gamma,b_\Gamma,\bar\gamma,L,\gamma$) and of
\cref{lem:stab-margins}\ref{lem:stab-margins:outer},
\ref{lem:stab-margins:roof} (with some $d_0,c_1,\tau$). Let $\zeta\ge0$,
$\delta<d_0$, $\delta+\zeta<\tau$, and let $S\subset K$ have
$\max(\ell^{\mathrm b}_K(t,p),\ell^{\mathrm d}_K(t,p))\ge-\zeta$ for $p\in S$
and $t\in(0,\pi/2)$. Then every point of $S$ lies within
$\max(1,1/c_1)(\delta+\zeta)$ of a point of $U_0$.
\item\label{lem:stab-directed:from} Suppose that $U_0$ has the property of
\cref{lem:stab-margins}\ref{lem:stab-margins:balls} (with some $\kappa,r_0$).
Let $A\ge0$, $\eps>0$, $\rho=(4(A+1)/\kappa)\sqrt\eps\le r_0$ and
$2\delta\le A\sqrt\eps$, and let $S$ be a set with
$\abs{U\setminus S}\le2\eps$. Then every point of $U_0$ lies within $\rho$ of a
point of $S$.
\end{enumerate}
\end{lemma}

\begin{proof}
\ref{lem:stab-directed:to} Let $p\in S\setminus U_0$; the coordinates of
$\ell_K$ and $\ell_{K_0}$ differ by at most $\delta$. If $p\in\cN(K_0)$, let
$d=\gamma(p_x)-p_y$ and take $t$ from
\cref{lem:stab-margins}\ref{lem:stab-margins:roof}; were $c_1d>\delta+\zeta$,
both coordinates of $\ell_K(t,p)$ would be less than $-\zeta$, so
$d\le(\delta+\zeta)/c_1$, and $(p_x,\gamma(p_x))\in U_0$ is at distance $d$.
Otherwise $p\notin K_0$, and some $q\in K_0$ has $\norm{p-q}\le\delta$
(\cref{lem:stab-lower}); were $q\in\cN(K_0)$, then
$\ip p{u_t}\le\ip q{u_t}+\delta\le h_{K_0}(t)-d_0+\delta<h_{K_0}(t)$ for
$t\in[0,\pi]$, and $p\in K_0$ by \eqref{eq:cap-fan}. So $q\in U_0$.
\ref{lem:stab-directed:from} First, $U$ contains every $p'$ with
$\mathbb B(p',2\delta)\subset U_0$: $p'+\delta u_t\in K_0$ gives
$\ip{p'}{u_t}\le h_K(t)$ for $t\in[0,\pi]$, so $p'\in K$; and if both
coordinates of $\ell_K(t,p')$ were negative, then $p'-\delta u_t-\delta v_t$,
at distance $\sqrt2\delta$, would have both coordinates of $\ell_{K_0}$
negative and lie in $\cN(K_0)$. Now let $p\in U_0$ have no point of $S$ in
$\mathbb B(p,\rho)$, take $\mathbb B(z,\kappa\rho)\subset U_0\cap\mathbb
B(p,\rho)$, and let $\mu=\kappa\rho/2=2(A+1)\sqrt\eps\ge2\delta$. The square of
side $\mu$ centered at $z$ lies in $\mathbb B(z,\mu)$, so the disks of radius
$2\delta$ about its points lie in $\mathbb B(z,\kappa\rho)$; hence it lies in
$U\setminus S$, and $4(A+1)^2\eps=\mu^2\le\abs{U\setminus S}\le2\eps$, which
is false.
\end{proof}

\begin{proposition}[local recovery of the sofa]\label{prop:stab-sofa}
There are $C,c,\delta,\psi_0,\eps_0>0$, with $\delta\le1$ and $\eps_0\le1$,
such that for every $K\in\Kc_{\pi/2}$ with $h_K(\pi)=h_{\KG}(\pi)$ and
$\dH(K,\KG)\le\delta$, every $\omega\in[0,\pi/2]$ with
$\pi/2-\omega\le\psi_0$ and every measurable set $S$ that satisfies the
constraints of $(K,\omega)$ and has $\eps=\abs\Gs-\abs S\in(0,\eps_0)$, the
sets $S$ and $\Gs$ are $C\sqrt\eps$-close and $\pi/2-\omega\le\eps/c$.
\end{proposition}

\begin{proof}
Take $c,\delta_1,\psi_0,R$ from \cref{prop:stab-terminal}, $\delta_2$ from
\cref{prop:stab-local}, $\delta=\min(\delta_1,\delta_2)$, and
$d_0,c_1,\tau,\kappa,r_0$ from \cref{lem:stab-margins}. Let $B'=4R/c$,
$C_1=\max(1,1/c_1)(80+B')$, $C_2=644/\kappa$ and $C=\max(1,C_1,C_2)$, and
choose $\eps_0\le1$ such that $80\sqrt\eps<d_0$, $(80+B')\sqrt\eps<\tau$ and
$C_2\sqrt\eps<r_0$ for $\eps<\eps_0$. By \cref{prop:stab-terminal},
$\pi/2-\omega\le\eps/c$, $\abs\Gs-\cA(K)\le\eps$,
$\abs{(K\setminus\cN(K))\setminus S}\le2\eps$, and
\ref{prop:stab-terminal:hall} holds with $4R(\pi/2-\omega)\le B'\sqrt\eps$, as
$\eps\le1$. As $s_K=0$, \cref{prop:stab-local,lem:stab-lower} give
$\dH(K,\KG)\le80\sqrt\eps$. By \cref{lem:stab-roof,lem:stab-margins}, $\KG$
has the properties of \cref{lem:stab-directed}, with
$\KG\setminus\cN(\KG)=\Gs$. Part \ref{lem:stab-directed:to}, with
$80\sqrt\eps$ for $\delta$ and $\zeta=B'\sqrt\eps$, puts $S$ within
$C_1\sqrt\eps$ of $\Gs$; part \ref{lem:stab-directed:from}, with $A=160$, puts
$\Gs$ within $C_2\sqrt\eps$ of $S$.
\end{proof}

\begin{lemma}[the symmetric difference]\label{lem:stab-symdiff}
For every $C>0$ there are $C'>0$ and $\eps_0\in(0,1]$ such that every compact
set $S$ that is $C\sqrt\eps$-close to $\Gs$, where
$\eps=\abs\Gs-\abs S\in(0,\eps_0]$, satisfies
$\abs{S\triangle\Gs}\le C'\sqrt\eps$.
\end{lemma}

\begin{proof}
Fix a disk $\mathbb B(z,r)\subset\KG$. Then
$h_{\KG}+d\le(1+d/r)h_{\KG}-(d/r)\ip z{u_t}$, so $\KG+d\,\mathbb B(0,1)$ lies
in the image of $\KG$ under the homothety of ratio $1+d/r$ about $z$. So the
points within $d\le1$ of $\KG$ outside it have area at most
$((1+d/r)^2-1)\abs{\KG}\le A_0d$ for a constant $A_0$. A point
$p\in\cN(\KG)$ within $d$ of a point $q\in\Gs$ has
$\gamma(p_x)-p_y\le(L+1)d$: if $q_x\in[a_\Gamma,b_\Gamma]$, then
$q_y\ge\gamma(q_x)\ge\gamma(p_x)-Ld$, and otherwise $\gamma(p_x)\le Ld$, while
$q_y\ge0$. With $C\sqrt{\eps_0}\le1$, $d=C\sqrt\eps$ and $d'=(L+1)d$, the
points of $S\setminus\Gs$ in $\cN(\KG)$ lie in
$\{(x,y):x\in[a_\Gamma,b_\Gamma],\ \gamma(x)-2d'<y<\gamma(x)+d'\}$, of area
$3d'(b_\Gamma-a_\Gamma)$. So $\abs{S\setminus\Gs}\le Ad$ with
$A=A_0+3(L+1)(b_\Gamma-a_\Gamma)$, and, as $\eps\le1$,
$\abs{S\triangle\Gs}=\abs\Gs-\abs S+2\abs{S\setminus\Gs}\le(1+2AC)\sqrt\eps$.
\end{proof}

\subsection{Compactness and uniqueness; proof of
\texorpdfstring{\cref{thm:stability,thm:stability-angle}}{the stability
theorems}}\label{sec:stab-entry}

Let $d_\infty$ be the Hausdorff distance between nonempty compact sets for the
norm $\norm{(x,y)}_\infty=\max(\abs x,\abs y)$. As $\norm p\le2\norm p_\infty$,
sets with $d_\infty\le r$ are $2r$-close; and the nonempty compact subsets of a
compact set form a compact space for $d_\infty$. For a compact set $S$,
$\ell^{\mathrm b}_S$, $\ell^{\mathrm d}_S$ and $\ell_S$ are defined as for
caps, with $h_S$ in place of $h_K$.

\begin{lemma}[limits of normalized sofas]\label{lem:stab-limit}
\begin{enumerate}[label=(\alph*)]
\item\label{lem:stab-limit:move} Let $S$ be a compact connected set and
$\omega\ge0$ with $S\subset\R\times[0,1]$, $h_S(\pi/2)=1$,
$h_S(\omega)+h_S(\omega+\pi)\le1$ and
$\max(\ell^{\mathrm b}_S(t,p),\ell^{\mathrm d}_S(t,p))\ge0$ for $p\in S$ and
$t\in[0,\omega]$. Then $S$ is a moving sofa with rotation angle $\omega$.
\item\label{lem:stab-limit:box} A moving sofa $S$ with rotation angle
$\omega\ge\pi/4$, $h_S(\pi/2)=1$ and $h_S(\pi)=h_\Gs(\pi)$ lies in
$[-h_\Gs(\pi),-h_\Gs(\pi)+6]\times[0,1]$.
\item\label{lem:stab-limit:lim} Let $S_n$ be moving sofas with rotation
angles $\omega_n\in[\omega_0,\pi/2]$, $h_{S_n}(\pi/2)=1$ and
$h_{S_n}(\pi)=h_\Gs(\pi)$, and let $S$ be a nonempty compact set with
$d_\infty(S_n,S)\to0$ and $\omega_n\to\omega$. Then $S$ is a moving sofa with
rotation angle $\omega\in[\omega_0,\pi/2]$, $h_S(\pi/2)=1$ and
$h_S(\pi)=h_\Gs(\pi)$; and if $\abs{S_n}\to M$, then $M\le\abs S$.
\end{enumerate}
\end{lemma}

\begin{proof}
\ref{lem:stab-limit:move} Let $\theta(s)=-s\omega$ and
$c(s)=-R_{-s\omega}\xx_S(s\omega)$. Then
$\Phi_s(p)=R_{-s\omega}(p-\xx_S(s\omega))$ has the coordinates
$\ell_S(s\omega,p)$, both at most $1$ and one nonnegative, so
$\Phi_s(S)\subset L$. At $s=0$ they are $(p_x-h_S(0)+1,p_y)$, so
$\Phi_0(S)\subset H_L$; at $s=1$, $\ell^{\mathrm b}_S(\omega,p)\in[0,1]$, as
$\ip p{u_\omega}\ge-h_S(\omega+\pi)\ge h_S(\omega)-1$, so
$\Phi_1(S)\subset V_L$.
\ref{lem:stab-limit:box} By \cref{lem:stab-sofacap}, $S\subset\R\times[0,1]$
and $S\subset L_S(\pi/4)$. As $S$ is connected and
$\ell^{\mathrm b}_S(\pi/4,\cdot)-\ell^{\mathrm d}_S(\pi/4,\cdot)$ takes both
signs on $S$ (where $\ip p{u_{\pi/4}}$ and where $\ip p{v_{\pi/4}}$ is
largest), some $z\in S$ has $\ell^{\mathrm b}_S=\ell^{\mathrm d}_S\ge0$ at
$\pi/4$, so $(h_S(\pi/4)+h_S(3\pi/4)-2)/\sqrt2=\xx_S(\pi/4)_y\le z_y\le1$.
For $p,q\in S$, as $p_y,q_y\ge0$,
$p_x-q_x\le\sqrt2(h_S(\pi/4)+h_S(3\pi/4))\le6$; and $\min_Sp_x=-h_\Gs(\pi)$.
\ref{lem:stab-limit:lim} Support values converge along converging
directions, as
$\abs{h_{S_n}(t_n)-h_S(t)}\le2d_\infty(S_n,S)+\abs{h_S(t_n)-h_S(t)}$
(\cref{lem:stab-lower}). So $h_S(\pi/2)=1$, $h_S(\pi)=h_\Gs(\pi)$, and
$h_S(\omega)+h_S(\omega+\pi)\le1$ by \cref{lem:stab-sofacap}; and
$S\subset\R\times[0,1]$. $S$ is connected, as disjoint closed sets that cover
$S$ and meet it have disjoint open neighborhoods that would eventually cover
$S_n$ and meet it. For $p\in S$ and $t\in[0,\omega]$, take $p_n\in S_n$ with
$p_n\to p$ and $t_n=(t/\omega)\omega_n$; the inequality
$\max(\ell^{\mathrm b}_{S_n},\ell^{\mathrm d}_{S_n})(t_n,p_n)\ge0$
(\cref{lem:stab-sofacap}) passes to the limit, and \ref{lem:stab-limit:move}
applies. Finally, for $r>\abs S$, an open set $\mathcal O\supset S$ with
$\abs{\mathcal O}<r$ contains $S_n$ for large $n$, so $M\le r$.
\end{proof}

\begin{proposition}[sofas of small deficit are near
$\Gs$]\label{prop:stab-entry}
For all $\rho,\psi_0>0$ there is $\eps_0>0$ such that for every moving sofa
$S$ with rotation angle $\omega\in[\omega_0,\pi/2]$ and $\eps(S)<\eps_0$, the
sets $S^\natural$ and $\Gs$ are $\rho$-close, and $\pi/2-\omega<\psi_0$.
\end{proposition}

\begin{proof}
Otherwise there are moving sofas $S_n$ with rotation angles
$\omega_n\in[\omega_0,\pi/2]$ and $\eps(S_n)<1/(n+1)$ that violate the
conclusion. The translates $S_n^\natural$ are moving sofas with the same
rotation angles, lie in the box of
\cref{lem:stab-limit}\ref{lem:stab-limit:box}, as $\omega_0\ge\pi/4$, and have
$\abs{S_n^\natural}\to\abs\Gs$. By compactness, a subsequence has
$d_\infty(S^\natural_n,S)\to0$ and $\omega_n\to\omega$ for a nonempty compact
set $S$; by \cref{lem:stab-limit}, $S$ is a moving sofa with rotation angle
$\omega$, $h_S(\pi/2)=1$, $h_S(\pi)=h_\Gs(\pi)$ and $\abs S\ge\abs\Gs$, so
$\abs S=\abs\Gs$ (\cref{fact:optimal}). \Cref{cor:translate} gives
$S+w=\Gs$ for some $w$, and the support values at $\pi/2$ and $\pi$ give
$w=0$. As $\Gs$ is a moving sofa with rotation angle $\omega$,
$\ip{q-p}{u_\omega}\le1$ on $\Gs$ (\cref{lem:stab-sofacap}), so
$\omega=\pi/2$ by \cref{lem:gerver-width}\ref{lem:gerver-width:gt}. For large
$n$, $d_\infty(S_n^\natural,\Gs)<\rho/2$ and $\pi/2-\omega_n<\psi_0$, a
contradiction.
\end{proof}

\begin{proof}[Proof of \cref{thm:stability,thm:stability-angle}]
Take $C,c,\delta,\psi_0,\eps_1$ from \cref{prop:stab-sofa}, $C',\eps_2$ from
\cref{lem:stab-symdiff} for this $C$, $\eps_3$ from \cref{prop:stab-entry} for
$\rho=\delta$ and this $\psi_0$, and $\eps_0=\min(\eps_1,\eps_2,\eps_3)$. Let
$S$ be a moving sofa with rotation angle $\omega\in[\omega_0,\pi/2]$ and
$\eps=\eps(S)<\eps_0$. If $\eps=0$, then $S^\natural=\Gs$ and $\omega=\pi/2$,
as in the proof of \cref{prop:stab-entry}. Otherwise $N=S^\natural$ is a
moving sofa with rotation angle $\omega$, and by \cref{prop:stab-entry} $N$
and $\Gs$ are $\delta$-close and $\pi/2-\omega<\psi_0$. The set $K$ of
\cref{lem:stab-sofacap} for $N$, a cap in $\Kc_{\pi/2}$, has
$h_K(\pi)=h_{\KG}(\pi)$ (as $h_{\KG}=h_\Gs$ on $[0,\pi]$) and
$\dH(K,\KG)\le\delta$, and $N$ satisfies the constraints of $(K,\omega)$. So
\cref{prop:stab-sofa} shows that $N$ and $\Gs$ are $C\sqrt\eps$-close and
$\pi/2-\omega\le\eps/c$, and \cref{lem:stab-symdiff} gives
$\abs{N\triangle\Gs}\le C'\sqrt\eps$. This proves \cref{thm:stability-angle},
with $C''=1/c$. For \cref{thm:stability}, let
$\eps_0'=\min(\eps_0,\abs\Gs-11/5)>0$
(\cref{fact:gerver}\ref{fact:gerver:def}): a moving sofa with deficit less
than $\eps_0'$ has area more than $11/5$, hence a rotation angle in
$[\omega_0,\pi/2]$ (\cref{fact:angle}).
\end{proof}

\subsection{Punctured sofas; proof of
\texorpdfstring{\cref{thm:stability-exponent}}{the optimality of the
exponent}}\label{sec:stab-puncture}

\begin{lemma}[removing a disk]\label{lem:stab-puncture}
Let $X$ be a closed connected set, $p\in\R^2$ and $r>0$ with
$\mathbb B(p,r)\subset X$, and let $Y=X\setminus\mathbb B^\circ(p,r)$. Then
$Y$ is closed and connected, and $r$-close to $X$; $\abs X-\abs Y=\pi r^2$ if
$X$ is compact; and if $X$ is a moving sofa with rotation angle $\omega$, so is
$Y$.
\end{lemma}

\begin{proof}
The boundary circle of $\mathbb B^\circ(p,r)$ lies in $Y$ and is connected. If
closed sets $A,B$ cover $Y$ and $Y\cap A\cap B=\emptyset$, the circle lies in
one of them, say $A$; then $(Y\cap A)\cup\mathbb B(p,r)$ and $Y\cap B$ are
disjoint closed sets that cover $X$, so $Y\cap B=\emptyset$. A point
$x\in\mathbb B(p,r)$ lies within $r$ of the point of the circle on the ray
from $p$ through $x$. The disk $\mathbb B^\circ(p,r)$ has area $\pi r^2$, and a
movement of $X$ moves the closed connected subset $Y$.
\end{proof}

\begin{lemma}[rigid copies keep interior points]\label{lem:stab-rigid}
Let $X$ be a compact set and $p$ an interior point of $X$. There is $\eta>0$
such that $p\in g(X)$ for every rigid motion $g$ for which $X$ and $g(X)$ are
$d$-close for some $d<\eta$.
\end{lemma}

\begin{proof}
Otherwise there are $g_n(x)=R_{\vartheta_n}x+v_n$ with $X$ and $g_n(X)$
$d_n$-close, $d_n\to0$ and $p\notin g_n(X)$. The $v_n$ are bounded, so a
subsequence has
$(\cos\vartheta_n,\sin\vartheta_n,v_n)\to(\cos\vartheta,\sin\vartheta,v)$;
let $g(x)=R_\vartheta x+v$. For $x\in X$, $g_n^{-1}(x)$ lies within $d_n$ of
$X$, so $g^{-1}(x)\in X$. Hence $g^{-1}$ maps the interior of $X$ onto an open
subset of $X$ that contains $g^{-1}(p)$, and, as $g_n^{-1}(p)\to g^{-1}(p)$,
$p\in g_n(X)$ for large $n$, a contradiction.
\end{proof}

\begin{proof}[Proof of \cref{thm:stability-exponent}]
\ref{thm:stability-exponent:family} As $\Gs$ is the closure of its interior
(\cref{prop:regular}), it has an interior point $p$; let
$\mathbb B(p,R)\subset\Gs$, let $\eta$ be as in \cref{lem:stab-rigid} for
$X=\Gs$, and let $r_0=\min(R/2,\eta/3)$ and $r\in(0,r_0)$. As $\Gs$ is a
moving sofa (\cref{fact:gerver}\ref{fact:gerver:cap},
\cref{fact:monotone}\ref{fact:monotone:envelope}), \cref{lem:stab-puncture}
shows that $\Gs_r$ is a moving sofa, $r$-close to $\Gs$, with
$\eps(\Gs_r)=\pi r^2$. If $\Gs_r$ and $g(\Gs)$ are $d$-close with $d<r$, then
$\Gs$ and $g(\Gs)$ are $(r+d)$-close and $r+d<\eta$, so $p\in g(\Gs)$ by
\cref{lem:stab-rigid}; then $p$ lies within $d<r$ of a point of $\Gs_r$, which
is false. So $r\le d$, and $d_{\mathrm{rig}}(\Gs_r,\Gs)=r$, attained by the
identity. \ref{thm:stability-exponent:no}, \ref{thm:stability-exponent:rig}
As $2a-1>0$, some $r\in(0,r_0)$ has $\pi r^2<\eps_0$ and
$C\pi^ar^{2a-1}<1$, that is, $C(\pi r^2)^a<r$. Then $S'=\Gs_r$ has
$0<\eps(S')<\eps_0$ and $C\eps(S')^a<r=d_{\mathrm{rig}}(S',\Gs)$, which is
\ref{thm:stability-exponent:rig}. Item \ref{thm:stability-exponent:rig}
implies \ref{thm:stability-exponent:no}: if $S'$ and $g(\Gs)$ were
$C\eps(S')^a$-close for a rigid motion $g$, then
$d_{\mathrm{rig}}(S',\Gs)\le C\eps(S')^a$. So $S=S'$ serves for
\ref{thm:stability-exponent:no}.
\end{proof}
